\documentclass[11pt]{article}
\usepackage[T1]{fontenc}
\usepackage{amsmath,amssymb,amsthm,geometry}
\geometry{margin=.9in}\setlength{\emergencystretch}{2em}
\allowdisplaybreaks[2]
\newtheorem{theorem}{Theorem}\newtheorem{lemma}{Lemma}
\newcommand{\E}{\mathbb E}\newcommand{\dd}{\,d}
\title{Cap-free control of the entrance term in the cross response}
\author{E5 continuation: analytic partial result}
\date{4 November 2026}
\begin{document}\maketitle

\paragraph{Result.}
The small bound for the actual weak cross-output $h$ through half
response is not closed. The homogeneous entrance contribution can,
however, be controlled without a residual-cap hypothesis. An exact
scalar damping formula and a contractive three-regression evolution
leave one explicit signed convolution of weak-target covariances.
Its needed smallness remains unproved. No residual exit or failure
of the desired initialized bound is asserted.

\section{An interval with no residual stop}
Use the actual aligned balanced isotropic critical family
$S_0=cq^{1/\omega}$, fixed $c>0$,
$\rho\in[13/20,7/10]$, $\lambda=\rho^2$,
$\omega=(1-\lambda)/2$. The exact model and loss are
\[
 f(D)=\int W(U^TD)_+\dd l,\quad |U|=|W|,\quad
 \mathcal L=\tfrac14\sum_{p=1}^2\E_{P_p}|D-f(D)|^2,
\]
with $P_1=N(e_1,q^2I)$ and $P_2=N(\rho e_2,q^2I)$.
Put $S=\int|U|^2\dd l$, $Q_1=1+q^2$, $Q=\lambda+q^2$, and
\[
 m=\E_{P_1}D_1f_1/Q_1,\quad z=\E_{P_1}D_1f_2/Q_1,
 \quad h=\E_{P_2}D_2f_1/Q,\quad n=\E_{P_2}D_2f_2/Q.
\]
Their full physical parameter gradients are $g_m,g_z,g_h,g$.
Write $w=Q(1-n)$ and $\mathcal R=4\mathcal L-Q(1-n)^2$.

Fix a resident delay $b$ and take the \emph{reached} time
$t_b=t_h+b$. Define
\[
 T_{1/2}=\inf\{t\ge t_b:n(t)=1/2\},\qquad
 T_K=\min\{T_{1/2},t_b+K\ell\},\quad \ell=\log(1/q),
\]
for any fixed $K>0$, with the infimum of the empty set equal to
$+\infty$. There is no stop at a residual cap, a chart boundary,
or a cross-output threshold. Neither finiteness of $T_{1/2}$ nor
its occurrence before the deterministic cutoff is assumed.

The fixed-time HBD/BRW entrance and SPM's cross-regression estimate give
\begin{equation}\label{an04cfcrv1:entry}
 m(t_b)\to M(b),\quad z(t_b),h(t_b)=O_b(q),\quad n(t_b)\to0,
 \quad 4\mathcal L(t_b)\to[1-M(b)]^2+\lambda,
 \quad M(b)=(1+e^{-b})^{-1}.
\end{equation}
Here $S(t_b)\le C_b$ and the global mass identity implies
\begin{equation}\label{an04cfcrv1:mass}
 S(t_b+s)\le C_b+Cs\quad(s\ge0),\qquad
 \int_{t_b}^{T_K}\|\dot\Theta\|^2\dd t\le\mathcal L(t_b)\le C.
\end{equation}
All regression coefficients are bounded by loss and Gaussian
projection inequalities, independently of a small residual cap.

There is even an order-one strong-response lower bound before
half response. Choose $b$ so that $[1-M(b)]^2<1/100$ and then
$q$ small so $4\mathcal L(t_b)\le\lambda+1/50$.
For $t\le T_K$, $n\le1/2$, even if $n$ is nonmonotone. Hence
\begin{equation}\label{an04cfcrv1:coarse}
 0\le\mathcal R(t)\le\lambda+1/50-Q/4\le31/80,
 \qquad m(t)\ge1-\sqrt{31/80}>1/3.
\end{equation}
This is a consequence of decreasing loss, not a retained
small-cap hypothesis. It does not imply the required
$\mathcal R\le r_*^2$, where $r_*=\lambda/800$.

\section{Exact finite-noise reductions remain valid without the cap}
REC's within-cluster residual projections give, directly at the
current state,
\begin{equation}\label{an04cfcrv1:flow}
 \dot\Theta=\tfrac{Q_1}{2}(1-m)g_m-\tfrac{Q_1}{2}zg_z
              -\tfrac Q2hg_h+\tfrac w2g+N_{\rm sh},
 \qquad \|N_{\rm sh}\|\le Cq\sqrt S(1+S).
\end{equation}
The proof uses only the network's $S$-Lipschitz bound and
Gaussian projection, so it does not require $\mathcal R\le r_*^2$.
The shape energy $\chi$ and the regression vector
\[
 e=\begin{pmatrix}\sqrt{Q_1}(m-1)\\\sqrt{Q_1}z\\\sqrt Qh\end{pmatrix}
\]
satisfy the exact identity
\begin{equation}\label{an04cfcrv1:regression}
 \mathcal R=|e|^2+\chi,\qquad 0\le\chi\le Cq^2(1+S)^2.
\end{equation}

For $b_p=\E_{P_p}D_p(U^TD)_+/Q_p$ and
$v_p=\E_{P_p}D_pD\mathbf1_{U^TD>0}/Q_p$, the gradient of
output-$i$ regression on cluster $p$ has blocks $(W_iv_p,e_ib_p)$.
Uniform Gaussian moment estimates, including all gate directions,
therefore give
\begin{align}
 \langle g,g_m\rangle&=O(qS),&
 \langle g_z,g_h\rangle&=O(qS),\nonumber\\
 \langle g_m,g_h\rangle&=C_{12}+O(qS),&
 \langle g,g_z\rangle&=C_{12}+O(qS),\nonumber\\
 \langle g,g_h\rangle&=T_2,&
 \|g_i\|&\le C\sqrt S,\label{an04cfcrv1:grams}
\end{align}
where $C_{12}=\int(U_1)_+(U_2)_+\dd l\ge0$ and
$T_2=\int W_1W_2|v_2|^2\dd l$ is signed.
These follow from $b_p=(U_p)_++O(q|U|)$ and
$v_p=\theta_pe_p+O(q)$, $0\le\theta_p\le1$.
They use physical mass factors, not isolated-label growth rates.

\section{First approach: the scalar homogeneous entrance does not amplify}
Taking the $g_h$ inner product of \eqref{an04cfcrv1:flow} gives
\begin{equation}\label{an04cfcrv1:scalar}
 h'+\tfrac Q2\|g_h\|^2h
 =\tfrac12\{wT_2+Q_1(1-m)C_{12}\}+\varepsilon_h,
 \qquad |\varepsilon_h|\le CqS(1+S).
\end{equation}
Specifically
\[
 \varepsilon_h=\tfrac{Q_1}{2}(1-m)
      (\langle g_h,g_m\rangle-C_{12})
 -\tfrac{Q_1z}{2}\langle g_h,g_z\rangle
 +\langle g_h,N_{\rm sh}\rangle.
\]
No signed term was dropped: the estimates in
\eqref{an04cfcrv1:grams} explain the paid errors.
Set
\[
 D_h(T,t)=\exp\left[-\tfrac Q2\int_t^T\|g_h(u)\|^2\dd u\right]
 \in(0,1].
\]
Then, uniformly on $[t_b,T_K]$,
\begin{equation}\label{an04cfcrv1:scalarvoc}
 h(T)=D_h(T,t_b)h(t_b)
 +\tfrac12\int_{t_b}^TD_h(T,t)
       [wT_2+Q_1(1-m)C_{12}](t)\dd t+O_{b,K}(q(1+\ell)^3).
\end{equation}
The first term is $O_b(q)$ on the entire interval, independently
of its length. At literal initialization isotropy gives
$f_0(D)=S_0D/4$ and hence $h(0)=0$ exactly. The corresponding
scalar formula from time zero has no homogeneous term at all.

This does not prove that physical transverse perturbations cannot
grow. GTM split $T_2$ into $h+(T_2-h)$ and placed the first part
in a potentially positive scalar coefficient. Equation
\eqref{an04cfcrv1:scalar} leaves all of that feedback in the
self-consistent forcing $wT_2/2$. Its sign and filtered size remain
to be controlled. The strong-overlap forcing in
\eqref{an04cfcrv1:scalarvoc} is also unpaid by this scalar argument.

\section{Second approach: absorb all regression couplings into dissipation}
Define the three parameter vectors
$j_1=\sqrt{Q_1}g_m$, $j_2=\sqrt{Q_1}g_z$,
$j_3=\sqrt Qg_h$, and their Gram matrix
$\mathsf G_{ij}=\langle j_i,j_j\rangle$.
It is positive semidefinite at every actual state. Differentiating
$e$ along the exact flow yields
\begin{equation}\label{an04cfcrv1:vector}
 e'=-\tfrac12\mathsf G e+\tfrac w2 b_0+\varepsilon,
 \qquad
 b_0=\begin{pmatrix}0\\\sqrt{Q_1}C_{12}\\\sqrt Q T_2\end{pmatrix},
 \qquad |\varepsilon|\le CqS(1+S).
\end{equation}
Here the explicit error components are
\[
 \varepsilon_i=\langle j_i,N_{\rm sh}\rangle
       +\tfrac w2(\langle j_i,g\rangle-(b_0)_i).
\]
In particular strong-response relaxation, its coupling to $h$,
and the mixed $z,h$ terms are included in $\mathsf G$.
None is estimated as an external unsigned source.

Let $\mathsf U(T,t)$ be the fundamental matrix for
$v'=-\mathsf Gv/2$, with $\mathsf U(t,t)=I$.
For each solution,
$\frac{d}{dT}|v|^2=-v^T\mathsf Gv\le0$, so
\begin{equation}\label{an04cfcrv1:contraction}
 \|\mathsf U(T,t)\|_{\rm op}\le1\qquad(T\ge t).
\end{equation}
No commutation or uniform positive eigenvalue is needed.
This is a pathwise equation with coefficients evaluated along
the actual trajectory, not its linearized sensitivity to a changed
initialization. The source $b_0$ depends on that same trajectory.

\begin{theorem}[Cap-free entrance and error control]
\label{an04cfcrv1:main}
For each fixed $c,\rho,b,K$ and all small $q$, the actual trajectory
on $[t_b,T_K]$ satisfies exactly
\begin{equation}\label{an04cfcrv1:representation}
 e(T)=\mathsf U(T,t_b)e(t_b)+\mathcal F(T)+\mathcal Z(T),
\end{equation}
where
\[
 \mathcal F(T)=\tfrac12\int_{t_b}^T
          w(t)\mathsf U(T,t)b_0(t)\dd t,\qquad
 \mathcal Z(T)=\int_{t_b}^T\mathsf U(T,t)\varepsilon(t)\dd t.
\]
Uniformly over this interval,
\begin{equation}\label{an04cfcrv1:paid}
 |\mathsf U(T,t_b)e(t_b)|\le1-M(b)+o_b(1),\qquad
 |\mathcal Z(T)|\le C_{b,K,c}q(1+\ell)^3=o(1).
\end{equation}
Thus the homogeneous entrance term can be made arbitrarily small
by taking fixed $b$ large and then $q$ small, without any residual
stop and without knowing the half-response hitting time.
\end{theorem}
\begin{proof}
Variation of constants applied to \eqref{an04cfcrv1:vector} proves
\eqref{an04cfcrv1:representation}. Apply
\eqref{an04cfcrv1:contraction} to both the entrance and the error
integral. Equation \eqref{an04cfcrv1:entry} gives
$|e(t_b)|\to1-M(b)$. Integration of
$CqS(1+S)$ using \eqref{an04cfcrv1:mass} proves the second bound.
All constants were selected on fixed $b,K$ before small $q$.
\end{proof}

\section{The precise signed forcing still uncontrolled}
The third component is
\begin{equation}\label{an04cfcrv1:unpaid}
 \boxed{\quad
 \mathcal F_h(T)=\tfrac12\int_{t_b}^T w(t)
 \left[\sqrt{Q_1}\mathsf U_{32}(T,t)C_{12}(t)
       +\sqrt Q\mathsf U_{33}(T,t)T_2(t)\right]\dd t.\quad}
\end{equation}
The exact $h$ identity is
\[
 \sqrt Qh(T)=[\mathsf U(T,t_b)e(t_b)]_3
                        +\mathcal F_h(T)+\mathcal Z_3(T).
\]
For example, choose $b$ with $1-M(b)<r_*/16$ and then $q$
small so the entrance term is below $r_*/8$ and the error below
$r_*/8$. A concrete sufficient component bound still needed is
\begin{equation}\label{an04cfcrv1:target}
 \sup_{t_b\le T\le T_K}|\mathcal F_h(T)|\le r_*/4.
\end{equation}
It would give $Qh(T)^2\le r_*^2/4$. This allocation is sufficient,
not necessary, and is not a full residual or E5 theorem.
The exact desired bound can instead be imposed on the signed sum
of all three terms in the preceding identity; neither has been
proved independently here.

Contraction gives only the rough estimate
$|\mathcal F|\le C\int_{t_b}^TS\dd t=O_{b,K,c}((1+\ell)^2)$.
Decreasing loss and \eqref{an04cfcrv1:coarse} improve it to
\begin{equation}\label{an04cfcrv1:insufficient}
 |\mathcal F(T)|\le\sqrt{31/80}+1-M(b)+o(1).
\end{equation}
This order-one estimate does not establish
\eqref{an04cfcrv1:target}, whose right side is $\lambda/3200$.
It is an insufficient upper bound, not a lower bound or a proof
that the forcing is large.

The sign $C_{12}\ge0$ does not sign its contribution after
transport by $\mathsf U_{32}$. Positive semidefiniteness of
$\mathsf G$ proves norm contraction, not entrywise positivity of
the evolution. $T_2$ is itself a signed covariance. Both terms
in \eqref{an04cfcrv1:unpaid} therefore remain in the same signed
integral. GTM's covariance bound by $\sqrt{S\Delta}$ and its signed
gate-transport estimates do not supply its required smallness at
these times. In particular absorbing the strong-overlap term into
$\mathsf G$ accounts for it but does not remove the transmitted
$C_{12}$ contribution.

\paragraph{What is and is not closed.}
The homogeneous entrance term and all finite-noise errors in this
representation are paid without the residual cap. The physical
feedback capable of amplifying $h$ has been retained in the
explicit actual-path forcing, not proved to cancel. That forcing
is the exact remaining term. The scalar and coupled-regression
approaches stop here. Neither an actual violation nor the desired
small bound through half response has been proved.

No boundary face has been assumed inward: the interval ends only
at the target half response or a deterministic horizon. The
order-one bound \eqref{an04cfcrv1:coarse} follows directly from
loss, including negative or nonmonotone $n$. The Gaussian feature
estimates include every gate direction. Fixed-time entrance is
used only at fixed $b$; no COD constant is claimed at
$s=t_\eta-t_h$. A reached $T_{1/2}$ later than the specified
logarithmic horizon is not covered by the uniform error bound,
and no half-response time bound is inferred from RQT's earlier
residual-stopped theorem. Constants and small-noise thresholds may
depend on fixed $c,\rho,b,K$, with no uniformity in extreme $c$.
Even a bound for $h$ would leave the other residual components and
the learned probe-depth inequality to control. Previous increments,
foundations, checkpoints, synced files, and GitHub remain unchanged.
All proof steps are analytic; no numerical trajectory or empirical
premise has been used.
\end{document}
